\documentclass[11pt]{article}

\usepackage[margin=1in]{geometry}
\usepackage[T1]{fontenc}
\usepackage[utf8]{inputenc}
\usepackage{amsmath,amssymb,amsthm,mathtools}
\usepackage{enumitem}
\usepackage{xr-hyper} % external references into the companion appendix document
\usepackage{hyperref}

\hypersetup{
  colorlinks=true,
  linkcolor=black,
  citecolor=black,
  urlcolor=black
}

\newtheorem{definition}{Definition}
\newtheorem{lemma}{Lemma}
\newtheorem{theorem}{Theorem}

% Resolve Lemma A.1--A.15 and Appendix A references against the
% separately compiled companion document (reads its .aux file):
\externaldocument{orp_np_hardness_appendix}

\newcommand{\E}{\mathbb{E}}
\newcommand{\Prb}{\mathbb{P}}
\newcommand{\ind}{\mathbf{1}}
\newcommand{\KT}{\operatorname{KT}}
\newcommand{\fas}{\operatorname{fas}}
\newcommand{\rev}{\operatorname{rev}}
\newcommand{\indeg}{\operatorname{indeg}}
\newcommand{\pos}{\operatorname{pos}}
\newcommand{\outdeg}{\operatorname{outdeg}}

\title{NP-Hardness of the Ordering Recovery Problem\\
\large A Corrected and Simplified Reduction}
\author{}
\date{}

\begin{document}
\maketitle

\begin{abstract}
This note gives a complete deterministic polynomial-time reduction from the minimum feedback arc set problem in tournaments to the Ordering Recovery Problem.  The construction uses one designated event for each tournament vertex, together with filler events.  The fillers make the designated events exactly independent under the posterior while contributing only neutral pairwise terms to the Kendall-tau objective.  Tournament arc directions are encoded by a McGarvey-style profile of $n(n-1)$ linear orders, and the posterior expected Kendall-tau distance becomes an affine rescaling of the FAST objective.  Instances with $n\le 3$ vertices are handled by brute force, so the main construction may assume $n\ge 4$.
\end{abstract}

\section{Preliminaries}

We write $S_L$ for the symmetric group on $\{1,\ldots,L\}$.  A permutation $\pi \in S_L$ will be interpreted as a ranking: $\pi(i)$ is the position assigned to item $i$, and item $i$ precedes item $j$ when $\pi(i)<\pi(j)$.

\begin{definition}[Minimum Feedback Arc Set in Tournaments]
A tournament is a directed graph $T=(V,A)$ such that for every unordered pair $\{u,v\}\subseteq V$ exactly one of $(u,v)$ and $(v,u)$ is in $A$.  For a linear order $\pi$ of $V$, define
\begin{equation}
  \fas(T,\pi)=\left|\{(u,v)\in A: \pi(u)>\pi(v)\}\right| .
\end{equation}
The FAST decision problem asks, given a tournament $T$ and an integer $k\ge 0$, whether there exists a linear order $\pi$ of $V$ with $\fas(T,\pi)\le k$.  FAST is NP-complete \cite{Alon2006,Charbit2007}.
\end{definition}

\begin{definition}[Ordering Recovery Problem]
An ORP instance has $L$ events $e_1,\ldots,e_L$.  A hidden permutation $\sigma\in S_L$ is drawn uniformly at random, and event $e_i$ occupies slot $\sigma(i)$, giving it true timestamp
\begin{equation}
  t_i^*=\sigma(i)\Delta,
\end{equation}
where $\Delta>0$ is fixed.  The observation is
\begin{equation}
  \hat t_i=t_i^*+\eta_i,
\end{equation}
where the noise variables $\eta_i\sim F_i$ are mutually independent and independent of $\sigma$.  In this note the distributions $F_i$ are finite-support rational probability mass functions.

Given the observation vector $\hat t=(\hat t_1,\ldots,\hat t_L)$ and a candidate output order $\hat\pi\in S_L$, the loss is the posterior expected Kendall-tau distance
\begin{equation}
  \E\left[d_{\KT}(\hat\pi,\sigma)\mid \hat t\right],
\end{equation}
where
\begin{equation}
  d_{\KT}(\hat\pi,\sigma)
  =\sum_{i<j}
  \ind\!\big[(\sigma(i)-\sigma(j))(\hat\pi(i)-\hat\pi(j))<0\big].
\end{equation}
The ORP decision problem asks whether there exists $\hat\pi\in S_L$ with
\begin{equation}
  \E\left[d_{\KT}(\hat\pi,\sigma)\mid \hat t\right]\le K,
\end{equation}
where $K$ is a non-negative rational threshold.
\end{definition}

\section{Pairwise decomposition}

Throughout the note, routine algebraic manipulations and auxiliary combinatorial facts are not established inline.  Instead they are isolated as Lemmas~\ref{alem:onesum}--\ref{alem:gadget} in Appendix~\ref{app:algebra}, which is maintained and compiled as a separate self-contained document (\texttt{orp\_algebra\_appendix.tex}); it contains only elementary identities and inequalities about real numbers, finite sequences of real numbers, and explicitly defined position maps, with no problem-specific material.  The main text supplies the constructive and probabilistic content and invokes the appendix lemmas with explicit instantiations of their variables, so that each appendix lemma can be verified in isolation.

The following two elementary lemmas are the only facts about the posterior expected Kendall-tau distance used in the reduction.

\begin{lemma}[Pairwise decomposition]\label{lem:pairwise}
For any prior distribution on $\sigma\in S_L$ and any observation vector $\hat t$,
\begin{align}
  \E\left[d_{\KT}(\hat\pi,\sigma)\mid \hat t\right]
  =\sum_{i<j} &\Prb(\sigma(i)<\sigma(j)\mid \hat t)\,\ind[\hat\pi(i)>\hat\pi(j)] \notag\\
  &+\Prb(\sigma(i)>\sigma(j)\mid \hat t)\,\ind[\hat\pi(i)<\hat\pi(j)].
\end{align}
\end{lemma}

\begin{proof}
For each pair $i<j$, let
\begin{align}
  X_{ij}(\sigma,\hat\pi)
  ={}&\ind[\sigma(i)<\sigma(j)\text{ and }\hat\pi(i)>\hat\pi(j)] \notag\\
   &+\ind[\sigma(i)>\sigma(j)\text{ and }\hat\pi(i)<\hat\pi(j)].
\end{align}
Then $X_{ij}=1$ exactly when $\sigma$ and $\hat\pi$ disagree on the relative order of $i$ and $j$, and $X_{ij}=0$ otherwise.  Thus $d_{\KT}(\hat\pi,\sigma)=\sum_{i<j}X_{ij}(\sigma,\hat\pi)$.  Taking conditional expectations and using that $\hat\pi$ is fixed gives the formula.
\end{proof}

\begin{lemma}[Majority form of a pairwise contribution]\label{lem:majority}
Fix $i<j$ and write
\begin{equation}
  p_{ij}=\Prb(\sigma(i)<\sigma(j)\mid \hat t).
\end{equation}
Call the more likely relative order the posterior majority order, breaking ties arbitrarily.  The contribution of pair $\{i,j\}$ to the posterior expected Kendall-tau distance is
\begin{equation}
  |2p_{ij}-1|\,\ind[\hat\pi\text{ disagrees with the posterior majority order on }\{i,j\}]
  +\min(p_{ij},1-p_{ij}).
\end{equation}
In particular, if $p_{ij}=1/2$, the contribution is always $1/2$, independent of $\hat\pi$.
\end{lemma}

\begin{proof}
By the previous lemma, the pair contribution is
\begin{equation}
  \ell_{ij}(\hat\pi)=p_{ij}\ind[\hat\pi(i)>\hat\pi(j)]
  +(1-p_{ij})\ind[\hat\pi(i)<\hat\pi(j)].
\end{equation}
If $\hat\pi$ agrees with the posterior majority order, this equals $\min(p_{ij},1-p_{ij})$.  If $\hat\pi$ disagrees, it equals $\max(p_{ij},1-p_{ij})$.  By Lemma~\ref{alem:maxmin} applied with $a=p_{ij}$ and $b=1-p_{ij}$, for which $|a-b|=|p_{ij}-(1-p_{ij})|=|2p_{ij}-1|$,
\begin{equation}
  \max(p_{ij},1-p_{ij})=\min(p_{ij},1-p_{ij})+|2p_{ij}-1|,
\end{equation}
and the claimed expression follows.  At $p_{ij}=1/2$ the coefficient $|2p_{ij}-1|$ is zero and the constant term is $1/2$.
\end{proof}

\section{The corrected reduction}

We reduce FAST to ORP.  Let $(T,k)$ be a FAST instance, where $T=(V,A)$ is a tournament on
\begin{equation}
  V=\{1,\ldots,n\}
\end{equation}
and let
\begin{equation}
  m=\binom n2 .
\end{equation}
For $u\in V$, write
\begin{equation}
  \delta_u=\indeg(u)-\outdeg(u).
\end{equation}
Thus $|\delta_u|\le n-1$.

We assume throughout that $n\ge 4$.  This is without loss of generality.  A many-one reduction must be defined on every input, so on instances with $n\le 3$ the reduction decides FAST by brute force in constant time (there are only finitely many such tournaments) and outputs a fixed YES instance or a fixed NO instance of ORP accordingly.  Since this branch affects only finitely many inputs and runs in constant time, it changes neither the polynomial running time nor the correctness of the reduction, and NP-hardness is an asymptotic statement in any case.

\subsection{A McGarvey profile with exact pairwise counts and rank sums}

We first build a list of linear orders of $V$ whose pairwise comparison counts encode the tournament.  The construction follows McGarvey's classical argument \cite{McGarvey1953}, which shows that every tournament arises as the majority relation of some list of linear orders.  For orders $\rho_1,\ldots,\rho_R$ of $V$ and an ordered pair $u\ne v$, write
\begin{equation}
  c_{uv}=\left|\{b:\rho_b(u)<\rho_b(v)\}\right|
\end{equation}
for the number of orders ranking $u$ before $v$; note $c_{uv}+c_{vu}=R$.

The construction is factored into two parts.  The \emph{constructive} part is a two-order gadget attached to each arc (Definition~\ref{def:gadget}); its bijectivity and its three combinatorial properties are stated and verified in abstract form in the companion appendix (Lemma~\ref{alem:gadget}).  The \emph{algebraic} part consists of two summation identities (Lemmas~\ref{alem:onesum} and~\ref{alem:threesum}).  The profile lemma (Lemma~\ref{lem:profile}) then follows by instantiating the appendix lemmas with the quantities supplied by the gadget construction.

\begin{definition}[Arc gadget]\label{def:gadget}
For a linear order $O$ of $V$ and $w\in V$, write $\pos_O(w)\in\{1,\ldots,n\}$ for the position of $w$ in $O$.  Let $(u,v)\in A$ and fix an arbitrary enumeration $x_1,\ldots,x_{n-2}$ of $V\setminus\{u,v\}$.  The \emph{gadget pair} of the arc $(u,v)$ is the pair of linear orders
\begin{equation}
  O_1=(u,\,v,\,x_1,\ldots,x_{n-2}),\qquad
  O_2=(x_{n-2},\ldots,x_1,\,u,\,v),
\end{equation}
that is, the orders defined by the position formulas
\begin{equation}\label{eq:gadget-positions}
\begin{aligned}
  \pos_{O_1}(u)&=1, & \pos_{O_1}(v)&=2, & \pos_{O_1}(x_j)&=2+j &&(1\le j\le n-2),\\
  \pos_{O_2}(u)&=n-1, & \pos_{O_2}(v)&=n, & \pos_{O_2}(x_j)&=n-1-j &&(1\le j\le n-2).
\end{aligned}
\end{equation}
By Lemma~\ref{alem:gadget}, applied with the symbols $u,v$ and the chosen enumeration $x_1,\ldots,x_{n-2}$ of $V\setminus\{u,v\}$, each row of \eqref{eq:gadget-positions} is a bijection $V\to\{1,\ldots,n\}$, so $O_1$ and $O_2$ are indeed linear orders of $V$; the same lemma records the combinatorial properties of the pair $(O_1,O_2)$ used below.
\end{definition}

\begin{lemma}[Profile construction]\label{lem:profile}
Set $R=2m=n(n-1)$, an even integer.  Enumerate the arcs as $A=\{a_1,\ldots,a_m\}$ and define the \emph{profile} to be the sequence of linear orders
\begin{equation}
  \rho_1,\ldots,\rho_R
\end{equation}
of $V$ obtained by letting $(\rho_{2c-1},\rho_{2c})$ be the gadget pair of arc $a_c$ from Definition~\ref{def:gadget}, for $c=1,\ldots,m$; as before, $\rho_b(u)$ denotes the rank position of $u$ in order $\rho_b$.  The profile is constructible in polynomial time and satisfies:
\begin{enumerate}[label=(\roman*)]
\item for every ordered pair $u\ne v$, if $(u,v)\in A$ then
\begin{equation}
  c_{uv}=\frac R2+1,
\end{equation}
and if $(v,u)\in A$ then
\begin{equation}
  c_{uv}=\frac R2-1;
\end{equation}
\item for every vertex $u$,
\begin{equation}
  \sum_{b=1}^R \rho_b(u)=\frac{R(n+1)}2+\delta_u.
\end{equation}
\end{enumerate}
\end{lemma}

\begin{proof}
The profile is well defined: it consists of the gadget pairs of all $m$ arcs, in the order of the chosen enumeration of $A$.  Each of its $R=2m$ orders is written down directly from the position formulas \eqref{eq:gadget-positions} in time $O(n)$, so the profile is constructible in polynomial time.  Since $T$ is a tournament, distinct arcs have distinct underlying unordered endpoint pairs; in particular, for any unordered pair $\{u,v\}$ exactly one arc has endpoint pair $\{u,v\}$, and every other arc $a$ satisfies $\{u,v\}\ne\{\text{endpoints of }a\}$.

\emph{Part (i).}  Fix an ordered pair $u\ne v$.  For each arc $a\in A$ with gadget pair $(O_1^a,O_2^a)$, define
\begin{equation}
  \gamma_a \;=\; \ind[\pos_{O_1^a}(u)<\pos_{O_1^a}(v)]+\ind[\pos_{O_2^a}(u)<\pos_{O_2^a}(v)]\;\in\;\{0,1,2\}.
\end{equation}
Each of the $R$ profile orders belongs to exactly one gadget pair, so
\begin{equation}\label{eq:cuv-split}
  c_{uv}=\sum_{a\in A}\gamma_a .
\end{equation}
By Lemma~\ref{alem:gadget}\ref{g:split}, $\gamma_a=1$ for each of the $m-1$ arcs $a$ whose endpoint pair differs from $\{u,v\}$.  For the unique arc with endpoint pair $\{u,v\}$: if that arc is $(u,v)$, then $\gamma_{(u,v)}=2$ by Lemma~\ref{alem:gadget}\ref{g:own}; if it is $(v,u)$, then $\gamma_{(v,u)}=0$, again by Lemma~\ref{alem:gadget}\ref{g:own} applied to the arc $(v,u)$ (both gadget orders place $v$ before $u$, so neither indicator in $\gamma_{(v,u)}$ is $1$).  The sum \eqref{eq:cuv-split} therefore consists of $m$ terms of which exactly one equals $2$ (resp.\ $0$) and the remaining $m-1$ equal $1$.  By Lemma~\ref{alem:onesum} with $c=2$ (resp.\ $c=0$),
\begin{equation}
  c_{uv}=(m-1)+2=m+1=\frac R2+1
  \qquad\text{resp.}\qquad
  c_{uv}=(m-1)+0=m-1=\frac R2-1,
\end{equation}
using $R=2m$.

\emph{Part (ii).}  Fix $u\in V$.  For each arc $a\in A$, define
\begin{equation}
  s_a \;=\; \pos_{O_1^a}(u)+\pos_{O_2^a}(u).
\end{equation}
Again because each profile order belongs to exactly one gadget pair,
\begin{equation}\label{eq:ranksum-split}
  \sum_{b=1}^R\rho_b(u)=\sum_{a\in A}s_a .
\end{equation}
By Lemma~\ref{alem:gadget}\ref{g:ranksum},
\begin{equation}
  s_a=
  \begin{cases}
    n, & a=(u,w)\text{ for some }w \quad(\text{$u$ is the tail of }a),\\
    n+2, & a=(w,u)\text{ for some }w \quad(\text{$u$ is the head of }a),\\
    n+1, & u\text{ is not an endpoint of }a.
  \end{cases}
\end{equation}
There are exactly $\outdeg(u)$ arcs of the first kind, $\indeg(u)$ of the second, and $m-\outdeg(u)-\indeg(u)$ of the third.  By Lemma~\ref{alem:threesum} applied to the sequence $(s_a)_{a\in A}$ with base value $n$, $d^+=\outdeg(u)$, and $d^-=\indeg(u)$,
\begin{equation}
  \sum_{a\in A}s_a=m(n+1)+\bigl(\indeg(u)-\outdeg(u)\bigr)=\frac{R(n+1)}2+\delta_u,
\end{equation}
using $R=2m$ and the definition $\delta_u=\indeg(u)-\outdeg(u)$.  Combining with \eqref{eq:ranksum-split} completes the proof.
\end{proof}

\subsection{The ORP instance}

Set
\begin{equation}
  L=2nR,
  \qquad q=L-n=2nR-n.
\end{equation}
The ORP instance has $L$ events: one designated event $d_u$ for each vertex $u\in V$, and $q$ filler events $f_1,
\ldots,f_q$.  Set $\Delta=1$ and set every observation equal to zero:
\begin{equation}
  \hat t_i=0\qquad\text{for every event }i.
\end{equation}

View the $L$ slots as $R$ consecutive blocks of width $2n$.  Block $b$ consists of slots
\begin{equation}
  (b-1)2n+1,\ldots,b2n.
\end{equation}
For vertex $u$ and block $b$, define the consecutive slot pair
\begin{equation}
  a_{u,b}^-=(b-1)2n+2\rho_b(u)-1,
  \qquad
  a_{u,b}^+=(b-1)2n+2\rho_b(u).
\end{equation}
Because $\rho_b$ is a linear order, the pairs $\{a_{u,b}^-,a_{u,b}^+\}$ are disjoint over $u$ inside each block.

Define
\begin{equation}
  \alpha_u=\frac12-\frac{2\delta_u}{R}+\frac{\delta_u}{R^2}.
\end{equation}
Write the perturbation as
\begin{equation}
  t_u=\frac{2\delta_u}{R}-\frac{\delta_u}{R^2},
  \qquad\text{so that}\qquad
  \alpha_u=\frac12-t_u .
\end{equation}
Since $|\delta_u|\le n-1$ and $R=n(n-1)\ge 1$, Lemma~\ref{alem:perturb}(ii) applied with $x=\delta_u$ and $M=n-1>0$ gives
\begin{equation}
  |t_u|<\frac{2(n-1)}{R}=\frac{2(n-1)}{n(n-1)}=\frac2n .
\end{equation}
Since $n\ge 4$, we have $2/n\le 1/2$, hence $|t_u|<1/2$, and $\alpha_u=\tfrac12-t_u\in(0,1)$ by Lemma~\ref{alem:perturb}(iii).

The noise distribution of designated event $d_u$ is
\begin{equation}
  F_{d_u}(-a_{u,b}^-)=\frac{1-\alpha_u}{R},
  \qquad
  F_{d_u}(-a_{u,b}^+)=\frac{\alpha_u}{R}
  \quad\text{for }b=1,\ldots,R,
\end{equation}
with probability zero on all other values.  The distribution sums to one because each block contributes $(1-\alpha_u)/R+\alpha_u/R=1/R$.

Each filler event has the uniform noise distribution
\begin{equation}
  F_f(-s)=\frac1L\quad\text{for }s=1,
\ldots,L,
\end{equation}
with probability zero elsewhere.  Since every observation is zero, assigning an event to slot $s$ requires noise value $-s$.  Thus, conditioned on the observation vector, the posterior weight of a permutation $\sigma$ is proportional to
\begin{equation}
  \prod_i F_i(-\sigma(i)).
\end{equation}
The construction is therefore exactly a finite-support rational ORP instance.

\subsection{Exact independence of the designated events}

Although the events compete for slots (no two events may occupy the same slot), the uniform fillers make the slots of the designated events mutually independent under the posterior.

\begin{lemma}[Independence via fillers]
Conditioned on the observation vector, the slots of the designated events are mutually independent.  More precisely,
\begin{equation}
  \Prb(X_u=a_{u,b}^-\mid \hat t)=\frac{1-\alpha_u}{R},
  \qquad
  \Prb(X_u=a_{u,b}^+\mid \hat t)=\frac{\alpha_u}{R},
\end{equation}
where $X_u$ is the slot occupied by $d_u$.
\end{lemma}

\begin{proof}
Let
\begin{equation}
  S_u=\{a_{u,b}^-,a_{u,b}^+: b=1,
\ldots,R\}
\end{equation}
be the support of designated event $d_u$.  The sets $S_u$ are pairwise disjoint over $u\in V$.

Fix any placement $(s_u)_{u\in V}$ with $s_u\in S_u$.  This placement uses $n$ distinct slots.  The $q=L-n$ fillers can be assigned bijectively to the remaining $q$ slots in exactly $q!$ ways.  Every such filler completion has the same filler likelihood $(1/L)^q$, and the number $q!$ of completions is independent of the chosen designated placement.  The prior on $\sigma$ is uniform and therefore also contributes only a placement-independent factor.

Consequently, the marginal posterior weight of the designated placement $(s_u)_{u\in V}$ is proportional to
\begin{equation}
  \prod_{u\in V} F_{d_u}(-s_u).
\end{equation}
Normalizing over the Cartesian product $\prod_u S_u$ shows that the joint posterior distribution of the designated slots is the product of the individual distributions, i.e.\ the slots are mutually independent, and each slot's posterior probabilities are exactly the designated noise weights.
\end{proof}

\subsection{Designated-designated posterior probabilities}

Let $B_u$ be the block containing $X_u$.  By the previous lemma, the random variables $B_u$ are mutually independent and uniform on $\{1,\ldots,R\}$.

Fix distinct vertices $u,v$.  If $B_u<B_v$, then $X_u<X_v$ because every slot in an earlier block is smaller than every slot in a later block.  If $B_u>B_v$, then $X_u>X_v$.  If $B_u=B_v=b$, then the comparison is decided by the profile order $\rho_b$: the two slots assigned to a vertex are consecutive, and the pairs for distinct vertices do not interleave.  Since $B_u$ and $B_v$ are independent and uniform on $\{1,\ldots,R\}$, every outcome $(B_u,B_v)=(b,b')$ has probability $1/R^2$; the event $B_u<B_v$ comprises $R(R-1)/2$ such outcomes by Lemma~\ref{alem:triangular}, and the event $\{B_u=B_v=b:\rho_b(u)<\rho_b(v)\}$ comprises $c_{uv}$ outcomes by the definition of $c_{uv}$.  Hence
\begin{align}
  \Prb(X_u<X_v\mid \hat t)
  &=\Prb(B_u<B_v)+\Prb(B_u=B_v\text{ and }\rho_{B_u}(u)<\rho_{B_u}(v)) \notag\\
  &=\frac{R(R-1)/2+c_{uv}}{R^2}.
\end{align}
Substituting $c_{uv}=R/2+1$ when $(u,v)\in A$ and $c_{uv}=R/2-1$ when $(v,u)\in A$ (Lemma~\ref{lem:profile}(i)), and applying Lemma~\ref{alem:normalized} with $\varepsilon=+1$ and $\varepsilon=-1$ respectively,
\begin{equation}
  \Prb(X_u<X_v\mid \hat t)
  =
  \begin{cases}
    \displaystyle \frac12+\frac1{R^2}, & (u,v)\in A,\\[6pt]
    \displaystyle \frac12-\frac1{R^2}, & (v,u)\in A.
  \end{cases}
\end{equation}
Therefore each designated pair has posterior majority direction equal to the tournament arc direction.  Writing $\Prb(X_u<X_v\mid\hat t)=\tfrac12+s_{uv}$ with $s_{uv}=\pm 1/R^2$, Lemma~\ref{alem:centered} gives its decision weight in Lemma~\ref{lem:majority} as
\begin{equation}
  \left|2\Prb(X_u<X_v\mid \hat t)-1\right|=2|s_{uv}|=\frac2{R^2},
\end{equation}
and its constant term as
\begin{equation}
  \min\left(\Prb(X_u<X_v\mid \hat t),1-\Prb(X_u<X_v\mid \hat t)\right)
  =\frac12-|s_{uv}|
  =\frac12-\frac1{R^2}.
\end{equation}

\subsection{Filler pairs are exactly neutral}

The remaining task is to show that every pair involving a filler has posterior probability $1/2$ in each direction.  Filler-filler neutrality follows from symmetry: fillers have identical distributions and the posterior is invariant under relabeling them, so for distinct fillers $f$ and $f'$,
\begin{equation}
  \Prb(X_f<X_{f'}\mid \hat t)=\Prb(X_{f'}<X_f\mid \hat t).
\end{equation}
Because $\sigma$ is a permutation there are no ties, so these two probabilities sum to one; being equal and summing to one, each equals $1/2$ by Lemma~\ref{alem:common}(iii).

Now fix a designated event $d_u$ and a filler $f$.  All fillers have the same noise distribution, so relabeling them leaves the posterior unchanged; hence
\begin{equation}
  \theta_u:=\Prb(X_f<X_u\mid \hat t)
\end{equation}
does not depend on which filler $f$ is chosen.  We compute $\theta_u$ using the pointwise rank identity
\begin{equation}
  X_u=1+
  \sum_{v\in V\setminus\{u\}}\ind[X_v<X_u]
  +\sum_{r=1}^q\ind[X_{f_r}<X_u].
\end{equation}
Taking expectations gives
\begin{equation}
  \E[X_u\mid \hat t]
  =1+
  \sum_{v\ne u}\Prb(X_v<X_u\mid \hat t)
  +q\theta_u.
\end{equation}

First compute $\E[X_u\mid \hat t]$.  Conditional on $B_u=b$, the slot $X_u$ takes the value $a_{u,b}^-=(b-1)2n+2\rho_b(u)-1$ with probability $1-\alpha_u$ and the value $a_{u,b}^+=a_{u,b}^-+1$ with probability $\alpha_u$, so by Lemma~\ref{alem:twopoint} (with $a=a_{u,b}^-$ and $\alpha=\alpha_u$),
\begin{equation}
  \E[X_u\mid \hat t,\;B_u=b]=(b-1)2n+2\rho_b(u)-1+\alpha_u.
\end{equation}
The block $B_u$ is uniform on $\{1,\ldots,R\}$, and the rank sums satisfy $\sum_{b=1}^R\rho_b(u)=R(n+1)/2+\delta_u$ by Lemma~\ref{lem:profile}(ii).  Lemma~\ref{alem:blockavg}, applied with $r_b=\rho_b(u)$, $\delta=\delta_u$, and $\alpha=\alpha_u$, therefore gives
\begin{equation}
  \E[X_u\mid \hat t]
  =\frac1R\sum_{b=1}^R\left((b-1)2n+2\rho_b(u)-1+\alpha_u\right)
  =nR+\frac{2\delta_u}{R}+\alpha_u.
\end{equation}
Substituting the definition $\alpha_u=\tfrac12-\tfrac{2\delta_u}{R}+\tfrac{\delta_u}{R^2}$ and cancelling the linear terms via Lemma~\ref{alem:cancel} (with $\delta=\delta_u$),
\begin{equation}
  \E[X_u\mid \hat t]=nR+\frac12+\frac{\delta_u}{R^2}.
\end{equation}

Next compute the designated-designated terms.  If $(v,u)\in A$, then $v$ is an in-neighbor of $u$ and
\begin{equation}
  \Prb(X_v<X_u\mid \hat t)=\frac12+\frac1{R^2}.
\end{equation}
If $(u,v)\in A$, then
\begin{equation}
  \Prb(X_v<X_u\mid \hat t)=\frac12-\frac1{R^2}.
\end{equation}
Thus, since $T$ is a tournament, $\indeg(u)+\outdeg(u)=n-1$ and $\indeg(u)-\outdeg(u)=\delta_u$, and Lemma~\ref{alem:halfsplit} applied with $p=\indeg(u)$, $q=\outdeg(u)$, and $s=1/R^2$ gives
\begin{align}
  \sum_{v\ne u}\Prb(X_v<X_u\mid \hat t)
  &=\indeg(u)\left(\frac12+\frac1{R^2}\right)
    +\outdeg(u)\left(\frac12-\frac1{R^2}\right) \notag\\
  &=\frac{\indeg(u)+\outdeg(u)}{2}+\frac{\indeg(u)-\outdeg(u)}{R^2}
  =\frac{n-1}{2}+\frac{\delta_u}{R^2}.
\end{align}
Substituting both computations into the rank identity,
\begin{equation}
  q\theta_u
  =\left(nR+\frac12+\frac{\delta_u}{R^2}\right)-1
    -\left(\frac{n-1}{2}+\frac{\delta_u}{R^2}\right).
\end{equation}
By Lemma~\ref{alem:common}(i) with $z=\delta_u/R^2$, the right-hand side equals $(2nR-n)/2=q/2$, using $q=2nR-n$.  Since $q=n(2R-1)>0$ (as $n\ge 4$ and $R\ge 1$), Lemma~\ref{alem:common}(ii) gives $\theta_u=1/2$.  Hence every filler-designated pair is exactly neutral:
\begin{equation}
  \Prb(X_f<X_u\mid \hat t)=\Prb(X_u<X_f\mid \hat t)=\frac12.
\end{equation}
By Lemma~\ref{lem:majority}, every pair involving a filler contributes exactly $1/2$ to the posterior expected Kendall-tau distance, independent of the candidate order $\hat\pi$.

\subsection{The objective identity}

Let
\begin{equation}
  C_0=m\left(\frac12-\frac1{R^2}\right)
  +\left(\binom L2-m\right)\frac12 .
\end{equation}
The $\binom L2=L(L-1)/2$ unordered pairs of events (Lemma~\ref{alem:triangular} with $R\coloneqq L$) partition into the $\binom n2=m$ designated-designated pairs and the $\binom L2-m$ pairs that involve at least one filler.  The first term of $C_0$ is the sum of the constant terms $\tfrac12-\tfrac1{R^2}$ over the $m$ designated-designated pairs; the second term is the fixed contribution $\tfrac12$ of each filler-involved pair.

Given any full event order $\hat\pi\in S_L$, let $\pi_{\hat\pi}$ be the induced order of the tournament vertices:
\begin{equation}
  \pi_{\hat\pi}(u)<\pi_{\hat\pi}(v)
  \quad\Longleftrightarrow\quad
  \hat\pi(d_u)<\hat\pi(d_v).
\end{equation}
This is a genuine linear order of $V$, because it is obtained by restricting a single permutation $\hat\pi$ to the designated events.

\begin{theorem}[Objective identity]\label{thm:objident}
For every $\hat\pi\in S_L$,
\begin{equation}
  \E\left[d_{\KT}(\hat\pi,\sigma)\mid \hat t\right]
  =C_0+\frac2{R^2}\fas(T,\pi_{\hat\pi}).
  \label{eq:objective-identity}
\end{equation}
\end{theorem}

\begin{proof}
By Lemma~\ref{lem:pairwise} and Lemma~\ref{lem:majority}, the posterior expected Kendall-tau distance is the sum over all $\binom L2$ pairs of a per-pair contribution of the form (decision weight)${}\times{}$(disagreement indicator) plus a constant.  The contributions are as follows.
\begin{itemize}
\item Each filler-involved pair has posterior probability exactly $1/2$ in each direction (Section~3.5), so by the final clause of Lemma~\ref{lem:majority} it contributes exactly $\tfrac12$, independent of $\hat\pi$.  The total over these pairs is $\left(\binom L2-m\right)\tfrac12$.
\item Each designated-designated pair $\{d_u,d_v\}$ contributes
\begin{equation}
  \left(\frac12-\frac1{R^2}\right)+\frac2{R^2}\,z_{uv},
  \qquad
  z_{uv}=\ind[\hat\pi\text{ disagrees with the posterior majority order on }\{d_u,d_v\}],
\end{equation}
by the designated-pair calculation of Section~3.4.  The posterior majority direction for $\{d_u,d_v\}$ is $d_u$ before $d_v$ exactly when $(u,v)\in A$, so $z_{uv}=1$ exactly when the corresponding tournament arc is backward under $\pi_{\hat\pi}$; hence
\begin{equation}
  \sum_{\{u,v\}} z_{uv}=\fas(T,\pi_{\hat\pi}),
\end{equation}
where the sum ranges over the $m$ unordered vertex pairs.  By Lemma~\ref{alem:affinesum} applied over this index set with $\kappa=\tfrac12-\tfrac1{R^2}$ and $w=\tfrac2{R^2}$, the total over designated-designated pairs is
\begin{equation}
  m\left(\frac12-\frac1{R^2}\right)+\frac2{R^2}\,\fas(T,\pi_{\hat\pi}).
\end{equation}
\end{itemize}
Adding the two totals and using the definition of $C_0$ yields \eqref{eq:objective-identity}.
\end{proof}

The theorem is the central point of the construction: the ORP objective is a constant plus a positive scalar multiple of the FAST objective, and the scalar is uniform over all arcs.

\subsection{Correctness of the reduction}

Set the ORP threshold to
\begin{equation}
  K=C_0+\frac{2k}{R^2}.
\end{equation}
We prove
\begin{equation}
  (T,k)\in\textsc{FAST}
  \quad\Longleftrightarrow\quad
  f(T,k)\in\textsc{ORP}.
\end{equation}

Suppose first that $(T,k)$ is a YES instance of FAST.  Then there is a linear order $\pi'$ of $V$ with $\fas(T,\pi')\le k$.  Choose any full event order $\hat\pi$ whose restriction to the designated events is $\pi'$; the fillers may be placed arbitrarily.  By \eqref{eq:objective-identity} and the right-to-left direction of Lemma~\ref{alem:threshold} (with $C=C_0$, $w=2/R^2>0$, $x=\fas(T,\pi')$),
\begin{equation}
  \E\left[d_{\KT}(\hat\pi,\sigma)\mid \hat t\right]
  =C_0+\frac2{R^2}\fas(T,\pi')
  \le C_0+\frac{2k}{R^2}=K.
\end{equation}
Thus the constructed ORP instance is a YES instance.

Conversely, suppose the constructed ORP instance is a YES instance.  Then there exists $\hat\pi\in S_L$ with
\begin{equation}
  \E\left[d_{\KT}(\hat\pi,\sigma)\mid \hat t\right]\le K.
\end{equation}
Let $\pi_{\hat\pi}$ be the induced order of the designated events.  Applying \eqref{eq:objective-identity},
\begin{equation}
  C_0+\frac2{R^2}\fas(T,\pi_{\hat\pi})
  \le C_0+\frac{2k}{R^2}.
\end{equation}
Since $2/R^2>0$, the left-to-right direction of Lemma~\ref{alem:threshold} (with $C=C_0$, $w=2/R^2$, $x=\fas(T,\pi_{\hat\pi})$) gives
\begin{equation}
  \fas(T,\pi_{\hat\pi})\le k.
\end{equation}
Hence $(T,k)$ is a YES instance of FAST.

Finally, the construction is polynomial time.  We have $R=2m=n(n-1)=O(n^2)$ and $L=2nR=2n^2(n-1)=O(n^3)$.  If the finite-support distributions are written explicitly, each designated distribution has support size $2R=O(n^2)$ and each filler distribution has support size $L=O(n^3)$, for a total explicit encoding size $O(L^2)=O(n^6)$.  All support points and rational probabilities have bit-length $O(\log n+\log k)$, except for the threshold numerator which also uses the input integer $k$.  Thus the map $(T,k)\mapsto f(T,k)$ is computable in deterministic polynomial time.  The brute-force branch for $n\le 3$ runs in constant time and is trivially correct, so the reduction is a total function that is polynomial-time computable and correct on all inputs.

\begin{theorem}
Let $(T,k)$ be a FAST instance on vertex set $V=\{1,\ldots,n\}$ with $n\ge 4$, and let $f(T,k)$ be the ORP instance constructed above from the profile $\rho_1,\ldots,\rho_R$ of Lemma~\ref{lem:profile}: $L=2nR$ events comprising one designated event $d_u$ per vertex $u\in V$, with noise distribution $F_{d_u}$ supported on the slot pairs $a^{\pm}_{u,b}$ with weights $(1-\alpha_u)/R$ and $\alpha_u/R$, and $q=L-n$ filler events with uniform noise, all observations equal to zero, together with the threshold $K=C_0+2k/R^2$.  Then $f$ is computable in deterministic polynomial time and
\begin{equation}
  (T,k)\in\mathrm{FAST}
  \quad\Longleftrightarrow\quad
  f(T,k)\in\mathrm{ORP}.
\end{equation}
In particular, the Ordering Recovery Problem is NP-hard under deterministic polynomial-time many-one reductions.
\end{theorem}

\begin{proof}
The equivalence is proved in Section~3.7 via the objective identity (Theorem~\ref{thm:objident}) and Lemma~\ref{alem:threshold}, and polynomial-time computability of the construction is verified at the end of Section~3.7.  For the NP-hardness conclusion, extend $f$ to the finitely many instances with $n\le 3$ by the constant-time brute-force branch described at the beginning of Section~3: decide $(T,k)$ by exhaustive search and output a fixed YES or NO instance of ORP accordingly.  The extension is correct by definition and affects neither the polynomial running time nor the correctness of $f$ on any input, so the extended $f$ is a total, deterministic polynomial-time many-one reduction from FAST to ORP.  Since FAST is NP-complete~\cite{Alon2006,Charbit2007}, ORP is NP-hard.
\end{proof}

\begin{thebibliography}{9}

\bibitem{Alon2006}
N. Alon.
\newblock Ranking tournaments.
\newblock \emph{SIAM Journal on Discrete Mathematics}, 20(1):137--142, 2006.

\bibitem{Charbit2007}
P. Charbit, S. Thomass\'e, and A. Yeo.
\newblock The minimum feedback arc set problem is NP-hard for tournaments.
\newblock \emph{Combinatorics, Probability and Computing}, 16(1):1--4, 2007.

\bibitem{McGarvey1953}
D. C. McGarvey.
\newblock A theorem on the construction of voting paradoxes.
\newblock \emph{Econometrica}, 21(4):608--610, 1953.

\end{thebibliography}

\end{document}
